
\section{Method}

\subsection{Overview}

Our CosformerAttention module has two computational paths:

\begin{enumerate}
\item \textbf{Encoder (bidirectional):} associative accumulation.
\item \textbf{Decoder (causal):} chunked prefix-sum.

\end{enumerate}
Both share the same feature map $\phi(\cdot) = \text{ELU} + 1$ (or softplus) 
and cosine/sine reweighting.

\begin{lstlisting}
                    Input X
                       |
        +--------------+--------------+
        |              |              |
                                       
    +------+       +------+      +------+
    |  Q   |       |  K   |      |  V   |
    | proj |       | proj |      | proj |
    +--+---+       +--+---+      +--+---+
       |              |              |
       |              |              |
                                     |
   +-----------------------+         |
   |   ( ) = ELU+1         |         |
   +---------+-------------+         |
             |                       |
                                     |
   +-----------------------+         |
   |  Cos/Sin reweighting  |         |
   +---------+-------------+         |
             |                       |
             +-----------+-----------+
             |           |           |
                                      
        +--------------------------------+
        | is_decoder?                    |
        +--+-------------------------+---+
           | No                      | Yes
           |                         |
                                      
    +--------------+          +--------------+
    | Associative  |          | Chunked      |
    | accumulation |          | prefix-sum   |
    +------+-------+          +------+-------+
           |                         |
           +------------+------------+
                        |
                         
                   Output Y
\end{lstlisting}

\subsection{Feature Map}

We support two feature maps:

\begin{equation}
\phi(x) = \begin{cases}
\text{elu}(x) + 1 & \text{(default)} \\
\text{softplus}(x) & \text{(alternative)}
\end{cases}
\end{equation}

Both are non-negative, deterministic, and cheap.

\textbf{Why not ReLU?} ReLU has zero derivative for negative inputs, which can 
cause dead neurons. ELU + 1 has a non-zero gradient everywhere.

\textbf{Why not RFF?} RFF is stochastic and requires more features for the same 
accuracy. ELU + 1 is deterministic and works well in practice.

\subsection{Positional Reweighting}

For sequence length $N$ and max position $M = \max(\text{max\_seq\_len}, N)$:

\begin{equation}
\theta_i = \frac{\pi i}{2 M}, \quad i = 0, 1, \dots, N-1
\end{equation}

The cosine and sine factors are:

\begin{equation}
c_i = \cos(\theta_i), \quad s_i = \sin(\theta_i)
\end{equation}

Applied to $Q, K$:

\begin{equation}
\tilde{Q} = Q \odot c, \quad \tilde{Q}' = Q \odot s, \quad
\tilde{K} = K \odot c, \quad \tilde{K}' = K \odot s
\end{equation}

Broadcasting: $c, s \in \mathbb{R}^{1 \times 1 \times N \times 1}$, 
$Q, K \in \mathbb{R}^{B \times H \times N \times d}$.

\subsection{Encoder Path: Associative Accumulation}

For bidirectional attention:

\begin{equation}
\begin{aligned}
KV_c &= \tilde{K}^\top V \in \mathbb{R}^{B \times H \times d \times d} \\
KV_s &= \tilde{K}'^\top V \in \mathbb{R}^{B \times H \times d \times d} \\
Z_c &= \sum_j \tilde{K}_j \in \mathbb{R}^{B \times H \times 1 \times d} \\
Z_s &= \sum_j \tilde{K}'_j \in \mathbb{R}^{B \times H \times 1 \times d}
\end{aligned}
\end{equation}

The output:

\begin{equation}
\text{num} = \tilde{Q} KV_c + \tilde{Q}' KV_s, \quad
\text{den} = \tilde{Q} Z_c^\top + \tilde{Q}' Z_s^\top + \epsilon
\end{equation}

\begin{equation}
\text{context} = \frac{\text{num}}{\text{den}}
\end{equation}

\textbf{Complexity.} Each matmul is $\mathcal{O}(L d^2)$ or $\mathcal{O}(d^2 L)$. 
Total: $\mathcal{O}(L d^2)$ time, $\mathcal{O}(d^2)$ memory for the state.

\subsection{Decoder Path: Chunked Prefix-Sum}

For causal attention:

\begin{enumerate}
\item Pad sequence to multiple of \texttt{chunk\_size = 32}.
\item Split into \texttt{num\_chunks = L / 32} chunks.
\item For each chunk:
\end{enumerate}
\begin{itemize}
\item Compute \textbf{intra-chunk} scores (with causal mask).
\item Compute \textbf{inter-chunk} contributions via running state $(S, Z)$.
\item Append output.
\item Update state.

\end{itemize}
\textbf{Intra-chunk:}

\begin{equation}
\text{num}_\text{intra} = \text{masked}(\tilde{Q}_c \tilde{K}_c^\top) V_c
\end{equation}

\textbf{Inter-chunk:}

\begin{equation}
\text{num}_\text{inter} = \tilde{Q}_c S_c + \tilde{Q}_c' S_s
\end{equation}

\textbf{State update:}

\begin{equation}
S_c \mathrel{+}= \tilde{K}_c^\top V_c, \quad
Z_c \mathrel{+}= \tilde{K}_c^\top \mathbf{1}
\end{equation}

The output for chunk $c$:

\begin{equation}
\text{context}_c = \frac{\text{num}_\text{intra} + \text{num}_\text{inter}}{\text{den}_\text{intra} + \text{den}_\text{inter} + \epsilon}
\end{equation}

\textbf{Complexity.} Outer Python loop over \texttt{num\_chunks} iterations (4 for $L = 128$). 
Each iteration does a small matmul. Total wall-clock: dominated by 4 kernel 
launches per chunk $$\times$ 2 chunks.

\subsection{Weight Transfer}

\subsubsection{Standard (DistilBERT)}

\begin{lstlisting}
q_w = src.q_lin.weight.data
k_w = src.k_lin.weight.data
v_w = src.v_lin.weight.data
o_w = src.out_lin.weight.data
\end{lstlisting}

\subsubsection{Nested (RoBERTa)}

\begin{lstlisting}
q_w = src.self.query.weight.data
k_w = src.self.key.weight.data
v_w = src.self.value.weight.data
o_w = src.output.dense.weight.data
\end{lstlisting}

\subsubsection{Fused (GPT-2)}

\begin{lstlisting}
c_w = src.c_attn.weight.data.t()
q_w, k_w, v_w = c_w.chunk(3, dim=0)
o_w = src.c_proj.weight.data.t()
\end{lstlisting}

\subsection{Parameter Analysis}

\begin{table}[H]
\centering
\small
\begin{tabular}{lll}
\toprule
\textbf{Component} & \textbf{Baseline} & \textbf{Cosformer} \\
\midrule
Q projection & $d^2$ & $d^2$ + LoRA \\
K projection & $d^2$ & $d^2$ + LoRA \\
V projection & $d^2$ & $d^2$ + LoRA \\
O projection & $d^2$ & $d^2$ + LoRA \\
\textbf{Per block} & $4d^2$ & $4d^2 + 4 \cdot r(d + d)$ \\
\bottomrule
\end{tabular}
\end{table}

\textbf{Example (DistilBERT, $r = 16$, $d = 768$):}

Per block: $4 \cdot 768\^{}2 + 4 \cdot 16 \cdot 1536 = 2{,}359{,}296 + 98{,}304 
= 2{,}457{,}600$.

$\Delta$\% per block: $+4.17\%$.

\textbf{Cosformer does not reduce parameter count.}

\subsection{Complexity Analysis}

\begin{table}[H]
\centering
\small
\begin{tabular}{lll}
\toprule
\textbf{Path} & \textbf{Time} & \textbf{Memory} \\
\midrule
Encoder & $\mathcal{O}(L d^2)$ & $\mathcal{O}(d^2)$ \\
Decoder & $\mathcal{O}(L d^2)$ & $\mathcal{O}(d^2)$ \\
Baseline MHSA & $\mathcal{O}(L^2 d)$ & $\mathcal{O}(L^2)$ \\
\bottomrule
\end{tabular}
\end{table}

\textbf{For $L = 128$, $d = 64$:}
\begin{itemize}
\item Cosformer: ~128 $$\times$ 64  = 524K operations per head.
\item Baseline: ~128  $$\times$ 64 = 1M operations per head.
\item \textbf{Cosformer has ~2$$\times$ fewer FLOPs.}

\end{itemize}
\textbf{Wall-clock:} Cosformer has overhead from:
\begin{itemize}
\item Kernel feature map (ELU)
\item Cos/sin reweighting
\item Chunked scan (decoder only)

\end{itemize}
\subsection{Algorithm}

\textbf{Algorithm 1:} CosformerAttention forward
\begin{lstlisting}
Input:  X $\in$ R^{B$$\times$L$$\times$d}, is_decoder, chunk_size C
Output: Y $\in$ R^{B$$\times$L$$\times$d}

1: Q $\leftarrow$ X W_Q;  K $\leftarrow$ X W_K;  V $\leftarrow$ X W_V
2: if RoPE: (Q, K) $\leftarrow$ Rotary(Q, K)
3: if GQA: (K, V) $\leftarrow$ repeat_kv(K, V)
4:
5: # Feature map
6: Q $\leftarrow$ ELU(Q) + 1
7: K $\leftarrow$ ELU(K) + 1
8:
9: # Positional reweighting
10: $\theta$ $\leftarrow$   * arange(L) / (2 * max_seq_len)
11: Q_c $\leftarrow$ Q   cos($\theta$); Q_s $\leftarrow$ Q   sin($\theta$)
12: K_c $\leftarrow$ K   cos($\theta$); K_s $\leftarrow$ K   sin($\theta$)
13:
14: if is_decoder:
15:     Y $\leftarrow$ chunked_prefix_sum(Q_c, Q_s, K_c, K_s, V, C)
16: else:
17:     KV_c $\leftarrow$ K_c^T @ V
18:     KV_s $\leftarrow$ K_s^T @ V
19:     Z_c $\leftarrow$ sum(K_c); Z_s $\leftarrow$ sum(K_s)
20:     num $\leftarrow$ Q_c @ KV_c + Q_s @ KV_s
21:     den $\leftarrow$ Q_c @ Z_c^T + Q_s @ Z_s^T
22:     Y $\leftarrow$ num / (den + $\epsilon$)
23:
24: if not is_nested_encoder:
25:     Y $\leftarrow$ out_proj(Y)
26: return Y
\end{lstlisting}

\subsection{Theoretical Notes}

\subsubsection{Kernel Approximation}

ELU + 1 is \textbf{not} a good approximation of $\exp(q^\top k)$. The 
approximation error is bounded by:

\begin{equation}
\left|\phi(q)^\top \phi(k) - e^{q^\top k}\right| \leq \epsilon_\text{ELU}
\end{equation}

where $\epsilon_\text{ELU}$ grows with $\|q\|$ and $\|k\|$.

\textbf{Why it works anyway.} The \textbf{softmax is scale-invariant} -- only the 
relative ordering of $q^\top k$ matters. ELU + 1 preserves the ordering 
approximately.

\subsubsection{Locality Recovery}

The Cosformer reweighting gives:

\begin{equation}
\tilde{q}_i^\top \tilde{k}_j + \tilde{q}_i'^\top \tilde{k}_j' = \phi(q_i)^\top \phi(k_j) \cdot \cos(\theta_i - \theta_j)
\end{equation}

The cosine term decays from 1 at $|i - j| = 0$ to 0 at $|i - j| = L$ 
(approximately). This restores locality.

\subsubsection{Causality}

The causal path preserves causality:

\begin{equation}
\frac{\partial Y_i}{\partial X_j} = 0 \text{ for } j > i
\end{equation}

because intra-chunk scores are masked and inter-chunk uses only 
previously-processed state.
\begin{lstlisting}

---
\end{lstlisting}